% Panduan Pengguna ISHAS — sumber LaTeX.
% Kompilasi: latexmk -pdf panduan-pengguna.tex   (atau: pdflatex 2x)
\documentclass[11pt,a4paper]{article}

\usepackage[T1]{fontenc}
\usepackage[utf8]{inputenc}
\usepackage[english,indonesian]{babel}
\usepackage{lmodern}
\usepackage{microtype}
\usepackage[a4paper,margin=2.4cm,footskip=1cm]{geometry}
\usepackage[table]{xcolor}
\usepackage{booktabs}
\usepackage{longtable}
\usepackage{array}
\usepackage{enumitem}
\usepackage{titlesec}
\usepackage{fancyhdr}
\usepackage[hidelinks]{hyperref}

\definecolor{ishas}{HTML}{9F1239}
\definecolor{ishaslight}{HTML}{FCE7F3}
\definecolor{headgrey}{HTML}{F3F4F6}

\titleformat{\section}{\color{ishas}\Large\bfseries}{\thesection}{0.6em}{}
\titleformat{\subsection}{\color{ishas}\large\bfseries}{\thesubsection}{0.6em}{}
\titleformat{\subsubsection}{\bfseries}{\thesubsubsection}{0.6em}{}

\pagestyle{fancy}
\fancyhf{}
\fancyhead[L]{\small\textcolor{ishas}{Panduan Pengguna ISHAS}}
\fancyhead[R]{\small\thepage}
\renewcommand{\headrulewidth}{0.4pt}

\setlist{topsep=2pt,itemsep=2pt,parsep=0pt}
\newcommand{\code}[1]{\texttt{#1}}

\title{\color{ishas}Panduan Pengguna ISHAS}
\author{Integrated Safety and Health Assessment System}
\date{\today}

\begin{document}

\begin{titlepage}
  \centering
  \vspace*{3cm}
  {\color{ishas}\Huge\bfseries Panduan Pengguna ISHAS\par}
  \vspace{0.6cm}
  {\large Integrated Safety and Health Assessment System\par}
  \vspace{0.4cm}
  {\large Sistem Penilaian K3L Pesantren\par}
  \vspace{2cm}
  {\large Versi dokumen: \today\par}
  \vfill
  \begin{minipage}{0.86\textwidth}
    \colorbox{ishaslight}{%
      \parbox{\dimexpr\linewidth-2\fboxsep}{%
        \textbf{Catatan status:} aplikasi masih \textbf{prototipe}. Semua angka,
        skor, kategori, dan isi penilaian adalah \textbf{data contoh}, bukan hasil
        penelitian atau rumus final. Sebagian data dapat hilang saat halaman
        di-refresh sesuai batas prototipe.%
      }%
    }
  \end{minipage}
  \vspace{1.2cm}
\end{titlepage}

\tableofcontents
\newpage

\section{Memulai}

\begin{itemize}
  \item Alamat aplikasi \code{/} (beranda) langsung menampilkan \textbf{dashboard publik} ---
        tidak ada halaman perkenalan.
  \item \textbf{Mode ngoding (dev):} halaman \code{/login} menampilkan
        \textbf{3 kartu akun contoh} --- cukup satu klik, tanpa kata sandi.
  \item \textbf{Mode rilis (prod):} halaman \code{/login} memakai
        \textbf{email + kata sandi}.
  \item Akun contoh (mode dev):
        \begin{itemize}
          \item Super Admin: \code{admin@ishas.demo}
          \item Validator: \code{validator@ishas.demo}
          \item Pesantren: \code{pesantren@ishas.demo}
        \end{itemize}
  \item Ingin berpindah peran? \textbf{Keluar} dulu, lalu masuk dengan akun lain.
        Tidak ada pemilih peran di dalam aplikasi.
\end{itemize}

\section{Publik / Pelapor (tanpa login)}

Siapa pun boleh membuka beranda dan halaman baca. Tidak perlu akun.

\subsection{Membaca dashboard publik (\code{/})}
\begin{itemize}
  \item Dashboard menampilkan ringkasan \textbf{semua pesantren terdaftar}
        (Aktif dan memiliki akun Pesantren aktif).
  \item Gunakan \textbf{pemilih pesantren} untuk memfilter seluruh angka/grafik ke
        satu pesantren. Pilihan ini juga tersimpan pada alamat
        (\code{?pesantren=PSN-0018}).
  \item Ada \textbf{pratinjau periode}; angka utama masih memakai periode berjalan
        (filter periode penuh menunggu keputusan D-04).
  \item Hanya laporan berstatus \textbf{Diterima/Terbit} yang tampil. Laporan
        \code{Menunggu validasi} dan \code{Ditolak} \textbf{tidak pernah} muncul di
        beranda.
\end{itemize}

\subsection{Mengirim laporan cepat (\code{/lapor})}
\begin{enumerate}
  \item Buka \code{/} atau \code{/lapor}, pilih pesantren pada pemilih.
  \item Isi satu formulir ringkas: nama pelapor, pesantren, area/lokasi, judul,
        deskripsi; kategori/aspek, usulan tingkat keparahan/prioritas/rekomendasi,
        foto bukti, dan kontak bersifat \textbf{opsional}.
  \item Tekan \textbf{Kirim laporan}. Anda menerima \textbf{nomor laporan}
        (\code{RPT-XXXX}) dan status \code{Menunggu validasi} --- laporan belum
        tampil publik sampai akun Pesantren menerimanya.
  \item Draf isian tersimpan di perangkat; jika keluar aplikasi, isian dapat lanjut.
\end{enumerate}

\noindent Catatan: laporan \textbf{tidak bisa} dikirim saat login sebagai Super
Admin/Validator --- keluar dulu untuk melapor sebagai publik. Akun Pesantren boleh
melapor (termasuk ke pesantren lain sebagai pelapor umum).

\subsection{Mengisi penilaian mandiri (\code{/penilaian-mandiri})}
\begin{enumerate}
  \item Daftarkan \textbf{nama penilai} dan pilih pesantren (kontak opsional).
  \item Jawab semua butir penilaian (6 bidang, 59 butir). Beberapa butir mewajibkan
        \textbf{bukti foto} dan/atau \textbf{lokasi}.
  \item \textbf{Draf} otomatis tersimpan; Anda bisa menutup halaman lalu melanjutkan.
  \item Setelah lengkap, tekan \textbf{Kirim penilaian}. Hasil \textbf{langsung
        terbit} tanpa validasi Pesantren dan skor/PDF langsung tampil publik.
        Kiriman yang sudah dikirim tidak dapat diubah.
\end{enumerate}

\subsection{Halaman baca publik}

\begin{longtable}{@{}>{\raggedright\arraybackslash}p{3.6cm}>{\raggedright\arraybackslash}p{11.2cm}@{}}
\toprule
\textbf{Halaman} & \textbf{Isi} \\
\midrule
\endhead
\code{/hasil} & Ringkasan nilai, nilai per bidang, dan daftar laporan PDF \\
\code{/peta-risiko} & Denah pesantren + titik temuan yang sudah sah, dengan filter \\
\code{/rekomendasi} & Daftar saran perbaikan: prioritas, penanggung jawab, kemajuan
                    (tenggat tidak ditampilkan publik) \\
\code{/tindak-lanjut} & Kemajuan perbaikan, status, nama penanggung jawab
                       (bukti penyelesaian tidak publik) \\
\code{/dokumen} & Kumpulan PDF tiap butir: yang \textbf{Umum} bisa dilihat/diunduh;
                 yang \textbf{Privat} terkunci \\
\code{/laporan/:id} & Laporan penilaian siap cetak: judul, nilai tetap, foto bukti,
                     tombol cetak/simpan PDF (tanpa temuan) \\
\code{/pesantren/\allowbreak PSN-XXXX} & Profil singkat satu pesantren (beranda terkunci ke
                            pesantren itu) \\
\bottomrule
\end{longtable}

\noindent Catatan privasi: nama pelapor, kontak, dan jawaban mentah tidak
ditampilkan publik. Kode pesantren yang tidak dikenal menampilkan pesan ``tidak
ditemukan'', bukan error.

\section{Pesantren (login)}

Akun admin lokal satu pondok. Setelah login, identitas akun tampil di kanan atas.

\subsection{Beranda (\code{/pesantren/dashboard})}
Ringkasan khusus pesantren Anda: jumlah antrean, status penanganan, dan kemajuan
tindak lanjut, beserta nama dan kode pesantren.

\subsection{Memeriksa laporan (\code{/pesantren/validasi-laporan})}
\begin{enumerate}
  \item Buka antrean laporan \textbf{lapor-cepat} milik pesantren Anda (penilaian
        mandiri tidak masuk antrean --- langsung terbit).
  \item Pilih laporan \code{Menunggu validasi}, lalu:
        \begin{itemize}
          \item \textbf{Terima} --- wajib menetapkan \textbf{Tingkat bahaya
                (severity)} dan \textbf{Prioritas (priority)} tanpa nilai bawaan,
                serta menulis rekomendasi final. Setelah diterima, laporan tampil di
                beranda publik dan \code{/tindak-lanjut}.
          \item \textbf{Tolak} --- wajib menulis \textbf{alasan minimal 10 karakter}.
                Laporan ditolak tidak tampil publik.
        \end{itemize}
  \item Anda \textbf{tidak boleh} mengubah isi laporan pelapor
        (deskripsi/jawaban/bukti); yang boleh diisi hanya keputusan, catatan
        validasi, atau alasan tolak.
\end{enumerate}

\subsection{Mengelola status penanganan (\code{/pesantren/tindak-lanjut})}
\begin{itemize}
  \item Alur status: \code{Pending} $\rightarrow$ \code{Proses} $\rightarrow$
        \code{Completed} (dan bisa dikembalikan dengan alasan). Laporan
        \code{Completed} dapat \textbf{diarsipkan} dengan alasan.
  \item Buat rencana tindakan: \textbf{PIC}, \textbf{tenggat}, dan catatan
        $\rightarrow$ status menjadi \code{Proses}.
  \item Perbarui \textbf{progres} (pilihan 0/25/50/75/100), unggah \textbf{bukti
        penyelesaian}, lalu ajukan verifikasi.
  \item Bisa \textbf{membatalkan perbaikan} dengan alasan minimal 10 karakter.
\end{itemize}

\subsection{Lokasi \& denah (\code{/pesantren/lokasi})}
Kelola gedung, area, dan \textbf{denah} (campus plan): unggah gambar lalu terbitkan.
Denah aktif dipakai untuk titik pelaporan dan peta risiko publik.

\subsection{Laporan \& hasil penilaian mandiri}
\begin{itemize}
  \item \code{/pesantren/laporan} --- daftar laporan cepat milik pesantren Anda.
  \item \code{/pesantren/hasil-penilaian-mandiri} --- daftar + rincian hasil
        penilaian mandiri milik pesantren Anda (hanya status \code{Terbit}).
\end{itemize}

\section{Validator (login)}

Mengelola instrumen/penilaian. Tidak memvalidasi laporan (itu tugas akun Pesantren).
\begin{itemize}
  \item \code{/validator/dashboard} --- ringkasan.
  \item \code{/validator/instrumen} --- \textbf{Bank instrumen}: tambah/ubah butir,
        tipe jawaban, pilihan + bobot, dan pengali. Perubahan langsung dipakai
        penilaian mandiri.
  \item \code{/validator/dokumen-instrumen} --- pustaka PDF per butir (Umum/Privat).
  \item \code{/validator/scoring} --- meninjau nilai akhir yang tersimpan; tautan
        ke laporan PDF.
  \item \code{/validator/validasi-publikasi} --- daftar periksa \textbf{5 syarat}
        sebelum hasil tampil ke umum.
  \item \code{/validator/data-penelitian} --- data penelitian, unduh CSV/JSON, dan
        impor bertahap (periksa $\rightarrow$ pratinjau $\rightarrow$ terapkan).
  \item \code{/validator/sam-isafe} (dan sub-halamannya) --- modul \textbf{SAM-iSAFE}:
        riwayat pengamatan + bank pertanyaan, formulir pengamatan baru (jawaban
        0/1/2), rincian skor per kategori, grafik perkembangan, riwayat perubahan,
        dan tindak lanjut temuan.
\end{itemize}

\noindent Catatan: pengisian SAM-iSAFE dapat dilanjutkan setelah halaman
di-refresh; tombol \textbf{Hapus draft} mengosongkan isian yang berjalan.

\section{Super Admin (login)}

Operator sistem pusat. Satu-satunya peran yang membuat pesantren dan akun.
\begin{itemize}
  \item \code{/admin/dashboard} --- ringkasan sistem.
  \item \code{/admin/pesantren} --- direktori pesantren: tambah, ubah status
        \code{Persiapan} $\rightarrow$ \code{Aktif} $\rightarrow$ \code{Nonaktif}.
        Hanya pesantren \code{Aktif} yang punya akun Pesantren aktif yang muncul di
        pemilih publik.
  \item \code{/admin/pengguna} --- kelola akun: buat (nama, email, sandi, peran,
        scope), ubah data, \textbf{reset sandi}, aktif/nonaktif, dan hapus (dengan
        konfirmasi).
  \item \code{/admin/hak-akses} --- tabel hak akses tiap peran (hanya baca).
  \item \code{/admin/audit-log} --- catatan aktivitas seluruh sistem (hanya baca).
  \item \code{/admin/pengaturan} --- \textbf{Reset data demo} dan pemindahan aset
        dari perangkat.
\end{itemize}

\section{Akun \& sesi}
\begin{itemize}
  \item \textbf{Akun aktif} selalu tampil di kanan atas (nama + peran + inisial).
  \item \textbf{Ganti kata sandi} (mode rilis): klik nama akun di kanan atas
        $\rightarrow$ \emph{Ganti kata sandi} $\rightarrow$ isi sandi lama + sandi
        baru (minimal 8 karakter) + konfirmasi. Setelah berhasil, sesi di perangkat
        lain keluar; perangkat ini tetap aktif.
  \item \textbf{Keluar}: klik nama akun $\rightarrow$ \emph{Keluar dari akun}.
        Anda diarahkan ke beranda publik (\code{/}).
  \item Refresh tidak mengeluarkan Anda dari sesi dan tidak menghapus draf isian.
\end{itemize}

\section{Warna \& label status}

Status selalu dibedakan dengan \textbf{label teks + ikon}, bukan warna saja.

\begin{center}
\begin{tabular}{@{}>{\raggedright\arraybackslash}p{4.2cm}>{\raggedright\arraybackslash}p{6.6cm}@{}}
\toprule
\textbf{Status} & \textbf{Keterangan} \\
\midrule
\cellcolor{headgrey}Menunggu validasi & Netral (ikon jam) \\
\cellcolor{yellow!25}Pending & Kuning (ikon jam) \\
\cellcolor{blue!15}Proses & Biru (sedang berjalan) \\
\cellcolor{green!15}Completed / Terverifikasi & Hijau (ikon centang) \\
\cellcolor{headgrey}Ditolak & Netral (silang; hanya terlihat di antrean Pesantren) \\
\cellcolor{headgrey}Dibatalkan & Netral (silang; tampil publik beserta alasannya) \\
\bottomrule
\end{tabular}
\end{center}

\noindent Tingkat bahaya/prioritas: \textcolor{red!70!black}{\textbf{Tinggi}},
\textcolor{orange!80!black}{\textbf{Sedang}},
\textcolor{green!50!black}{\textbf{Rendah}}, dan \textbf{Belum ditentukan} (netral).

\section{Batasan prototipe}
\begin{itemize}
  \item Semua angka/skor/kategori adalah \textbf{ilustrasi}, bukan rumus ilmiah final.
  \item Beberapa data (mis. foto bukti) hanya tersedia di perangkat pengunggah;
        setelah berpindah perangkat, foto bisa tidak tampil.
  \item Biaya, domain, dan hosting pada nota terpisah tidak memengaruhi cara pakai
        aplikasi ini.
\end{itemize}

\end{document}
